\section{Introduction}

Academic writing is a complex document engineering process. In disciplines that
use \LaTeX{}, a manuscript is not merely a linear text document: it is a
project composed of source files, bibliographic databases, figures, style
templates, experimental code, compiler logs, reviewer feedback, and final PDF
artifacts. Researchers must coordinate editing, citation management, compilation,
revision tracking, and collaboration across these interconnected components.

The emergence of large language models (LLMs)~\cite{openai2023gpt4,
anthropic2024claude, openai2024gpt4o} has introduced new possibilities for
AI-assisted writing. Tools such as ChatGPT, Claude, Gemini, and specialized
academic assistants can generate, paraphrase, summarize, and critique scholarly
text. However, general-purpose chat interfaces treat writing as an isolated
text-generation task. They separate model output from the project context in
which scholarly artifacts must be edited, cited, compiled, and reviewed. This
separation creates three practical problems:

\begin{enumerate}
  \item \textbf{Context fragmentation.} Researchers must repeatedly copy
  manuscript sections into the assistant, increasing friction and the risk of
  stale or incomplete prompts. When the manuscript changes, chat context
  diverges from source truth.

  \item \textbf{Unaudited edits.} When assistants can directly modify files, or
  when changes are manually copied into a source tree, it becomes difficult to
  track what was changed, by whom, and why. This undermines the audit trail
  essential for scholarly integrity.

  \item \textbf{Unverifiable artifacts.} Generated text may introduce
  hallucinated references, missing citation keys, or \LaTeX{} changes that fail
  to compile. Liang et al.~\cite{liang2024llmwriting} documented widespread
  use of LLMs in scientific manuscripts and highlighted significant risks of
  hallucinated references, underscoring the verification gap. These
  are document engineering failures that text-generation quality metrics do not
  capture.
\end{enumerate}

Existing solutions address parts of this problem---Overleaf~\cite{overleaf}
provides collaborative \LaTeX{} editing with recently introduced AI
features~\cite{overleaf2024ai}; citation managers such as Zotero and Mendeley
handle references; and AI chat interfaces generate text. Yet no existing
system integrates these capabilities into a unified, controllable, local-first
workspace designed for the full manuscript lifecycle.

We present \software{}, a local-first, human-in-the-loop software platform
that addresses this gap by treating academic writing as a software-backed
document engineering workflow. The system is designed and implemented
following practical software engineering principles: a modular monorepo
architecture, comprehensive automated testing across 32 test files, secure
configuration management with masked credentials, and careful attention to
real-world deployment concerns such as API rate limiting, error recovery,
and dependency management. The system represents each paper as a
filesystem-backed project and provides an integrated workspace for editing,
preview, compilation, citation verification, AI assistance, human checkpoints,
and workflow orchestration. AI interaction is separated into three
permission-aware modes to give researchers fine-grained control over how and
when model output enters their manuscript.

This paper makes the following contributions:

\begin{itemize}
  \item \textbf{System architecture.} We describe the design and architecture
  of \software{}, a local-first platform integrating AI-assisted writing with
  project-level document management, citation verification against multiple
  scholarly databases, \LaTeX{} compilation, and tmux-based terminal
  management. The implementation comprises approximately 23,100 lines of application
  source code across a monorepo with React/Vite frontend and Fastify/Node.js backend.

  \item \textbf{Permission-aware interaction model.} We present a
  three-mode AI interaction model (Chat, Agent, Tools) that separates read-only
  discussion, reviewable edit proposals with inline diff display, and
  controlled tool execution. This design emerged from iterative user feedback
  during development and is designed to reduce anxiety about unintended
  manuscript changes and promote exploratory AI use.

  \item \textbf{Typed workflow pipelines.} We report on Pipeline 2.0, a
  system of typed pipeline stages (AI, Compute, Human, Citation, Compile) with
  human checkpoints as first-class execution gates, enabling auditable
  multi-stage writing processes.

  \item \textbf{Practical experience.} We share lessons learned from building
  and using the system across multiple development iterations,
  including design tradeoffs, testing strategies, the challenges of LLM
  integration, and validation through a 32-file automated test suite.
\end{itemize}

The remainder of this paper is organized as follows. Section~\ref{sec:related}
reviews related work. Section~\ref{sec:design} presents the system design and
architecture. Section~\ref{sec:implementation} describes implementation
details. Section~\ref{sec:evaluation} reports evaluation results and practical
experience. Section~\ref{sec:discussion} discusses lessons learned.
Section~\ref{sec:limitations} outlines limitations, and Section~\ref{sec:conclusion}
concludes.
